%% APPENDIX C -- the line mask.  Carries app:linemask (04_method.tex,
%% 05_results.tex, app_A, app_J), eq:mask and tab:masklines.
%%
%% ★ THIS APPENDIX IS THE PAPER'S ONLY DEFINITION OF THE ATTRIBUTION RULE.
%% It owns the species list, the half-width and the FRAME.  The rule is the
%% star's own rest frame; the local standard of rest is evaluated as a check
%% and changes no disposition (masklsr_v520.py, L10).
%% r30 (write-apx): Table C1, the rule as restated (list completed after the
%% crossings were examined), the definition as one equation, the tuning
%% caveat once, the LSR check and the half-width with its cost.
%% r35 (W3, m5): the selection rule as the code applies it (every hyperfine
%% component its own entry, no strength cut): \FrNTrans entries = \McNRow transitions (N4, maskso S10/S11).  That the
%% probabilities are conditional on this adopted mask is stated once, in
%% Sec. 4.3.1, and not here.
\section{The line mask}
\label{app:linemask}

A crossing that falls on a known molecular transition has an astrophysical
explanation to hand, and the mask decides which crossings those are.

The species list was assembled by hand. It holds the bright
millimetre disc and cloud tracers the pipeline checked from its first search,
with [C\,\textsc{i}] and hydrogen recombination added during the search and SO
after a crossing was found on one of its lines. It was completed at the query
below, after the crossings had been examined, and it omits other species seen
in discs and clouds, such as CH$_3$OH, DCO$^{+}$ and N$_2$H$^{+}$. The mask
contains every transition of these \FrNSpec{} species inside the
\McNIsland{} searched intervals with an upper-state energy below
\McEuMax\,K, \FrNTrans{} catalogue entries in all
(Table~\ref{tab:masklines}). There is no strength cut, and every hyperfine
component the catalogues list is an entry of its own, so the \FrNTrans{}
entries are \McNRow{} transitions, of which \MkHfsNTrans, in
\MkHfsSpecList, are split into components. All come from one Splatalogue query \citep{Splatalogue2007} of CDMS and JPL
\citep{Pickett1998,CDMS2005} dated \MkListDate.

Attribution is decided in the star's own rest frame. A crossing's
topocentric frequency is carried to the barycentre with its own block's
pointing and mid-time, and then to the star with its systemic velocity,
giving $\nu_\star$. The crossing is attributed if, for any listed transition
of rest frequency $\nu_0$,
\begin{equation*}
|\Delta v_\star|\le v_{\rm h},\qquad
\Delta v_\star=c\,\frac{\nu_\star-\nu_0}{\nu_0},\qquad
v_{\rm h}=\MaskHalfKms\,{\rm km\,s^{-1}},
\end{equation*}
so that $\Delta v_\star$ is positive for a crossing above the transition in
frequency. The same test, in the same frame, decides whether a control
maximum is attributed (Appendix~\ref{app:falsealarm}). Of the
\CsNCrossQp{} quality-passing census crossings, \MkNCoinc{} meet it, and these
are the \MkAttr{} attributions; the hold-out is treated identically.

The windows were placed by their observers for their own lines, and that
limits what ``unattributed'' means. Of the Class~A windows, \TnCenPct{} per
cent are centred within $v_{\rm h}$ of a listed transition, against
\TnCtrlCenLo--\TnCtrlCenHi{} per cent for the same windows displaced, so a
crossing is unattributed against a list that these setups helped to shape.

As a check that changes no disposition, every crossing is also read in the
local standard of rest. Toward
\LsCloudStars{} the CO($1\to0$) survey of \citet{DameThaddeus2022,Dame2001}
shows emission at the crossing's own LSR velocity (within
\LsCloudDvMax\,km\,s$^{-1}$ of its peak, and at least \LsCloudTMin\,K), so
\LsNCloud{} crossings toward these stars lie at the velocity of
interstellar gas. Elsewhere the Planck CO
map \citep{PlanckCO2016} stays below
\LsPlanckMax\,K\,km\,s$^{-1}$.

Keplerian motion where these molecules survive sets the half-width,
$\lesssim$\EcKepVKms\,km\,s$^{-1}$ beyond \EcKepRAu\,au. It has a cost.
Over the union of searched sky frequencies the mask removes \MkMaskAGHz\,GHz
of the \CvUnionA\,GHz Class~A union, \MkMaskAPct{} per cent, and
\FrMaskLostGHz\,GHz of \FrMaskUnionGHz\,GHz over both classes, \FrMaskPct{}
per cent, although each limit is formed over its own window's channels
outside the mask in its own star's frame (\S\ref{sec:summary}).

\begin{table}
\centering
\footnotesize
\caption{The line mask by species: catalogue entries $N$ (each hyperfine
component one entry), transitions $n_{\rm t}$, and the frequencies and
upper-state energies they span.}
\label{tab:masklines}
%% GENERATED by maskso_v414.py -- do not hand-edit.
\begin{tabular}{@{}lrrr@{\,}c@{\,}rr@{}}
\hline
Species & $N$ & $n_{\rm t}$ & \multicolumn{3}{c}{$\nu_0$ (GHz)} & $E_{\rm u}$ (K) \\
\hline
H & 5 & 5 & 92.03 & -- & 231.90 & -- \\
SO & 15 & 15 & 100.03 & -- & 682.68 & 16--143 \\
H$_2$CO & 13 & 7 & 101.33 & -- & 491.97 & 10--106 \\
C$^{18}$O & 3 & 3 & 109.78 & -- & 329.33 & 5--32 \\
CN & 50 & 10 & 113.12 & -- & 680.30 & 5--114 \\
CO & 3 & 3 & 115.27 & -- & 345.80 & 6--33 \\
CS & 3 & 3 & 195.95 & -- & 342.88 & 24--66 \\
SiO & 2 & 2 & 217.10 & -- & 347.33 & 31--75 \\
$^{13}$CO & 8 & 2 & 220.40 & -- & 330.59 & 16--32 \\
HCN & 7 & 1 & 354.50 & -- & 354.51 & 43 \\
HCO$^{+}$ & 1 & 1 & 356.73 & -- & 356.73 & 43 \\
{[C\,\textsc{i}]} & 1 & 1 & 492.16 & -- & 492.16 & 24 \\
\hline
\end{tabular}%

\end{table}
